\section*{How to Read and Extend These Notes}
\addcontentsline{toc}{section}{How to Read and Extend These Notes}

These notes are intended to complement a lecture presented at the board. The lecture can emphasize the main physical argument, develop a representative derivation, and respond to questions. The written notes can preserve the complete structure: definitions, conventions, intermediate steps, calculations, checks, diagrams, and conclusions. Neither form should merely duplicate the other. Together they should allow the student to see both the main line of reasoning and the details needed for independent study.

A productive way to read a chapter is to separate three levels of understanding.

\begin{enumerate}
    \item \textbf{Conceptual level:} What physical or mathematical problem is being addressed? What new object or principle is introduced?
    \item \textbf{Derivational level:} How does the main equation follow from the assumptions? Which steps use algebra, calculus, symmetry, or boundary conditions?
    \item \textbf{Operational level:} Given a specific field, source, or transformation, can the method be applied to obtain a concrete result?
\end{enumerate}

The student should not try to memorize every displayed equation at the first reading. It is more useful to identify the role of each equation. Is it a definition, an identity, an equation of motion, a conservation law, a gauge choice, or a solution under specified boundary conditions? Once that role is understood, the formula is easier to reconstruct and use.

Calculations should be read actively. Before following a derivation, the reader should note the variables, free indices, dimensions, and expected form of the answer. After the derivation, the result should be checked in a simple case. For instance, a covariant equation can be expanded into components, a conservation law can be integrated over space, and a relativistic expression can be examined in a low-velocity limit.

The notes will be developed incrementally. Each chapter is divided into section files, and the substantive content will be written one section at a time. This is not only a repository-management choice. It supports mathematical and pedagogical quality. A single section can be checked for logical continuity, notation, sign conventions, examples, and links to earlier material before the next section is added.

The intended development cycle for future work is:

\begin{enumerate}
    \item select one section;
    \item identify its precise learning objective;
    \item write the motivation and definitions;
    \item develop the derivation or central argument;
    \item add one or more explicit examples;
    \item verify notation, equations, and limiting cases;
    \item compile and inspect the PDF;
    \item continue to the next section only after the current one is coherent.
\end{enumerate}

This section-by-section method prevents a chapter from becoming a long sequence of loosely connected formulas. It also makes it possible to revise the course structure without rewriting unrelated material. If a topic later requires more detail, its section can be expanded or divided. If a topic proves secondary, it can be shortened without disturbing the rest of the chapter.

The twelve chapters describe the intended logical route, but the final notes should remain responsive to the actual needs of students. Some ideas may require additional preliminary explanations. Some calculations may need a second example. A diagram may communicate a geometric relation more effectively than another paragraph. The structure should therefore be stable enough to guide the course, but flexible enough to improve as the material is taught and reviewed.

A successful use of these notes should lead to more than recognition. After studying a topic, the student should be able to explain the problem in words, reproduce the main derivation with stated assumptions, carry out a representative calculation, interpret the result, and recognize the most common mistakes. That standard will guide the construction of every subsequent section.
